\chapter{\ifenglish Background Knowledge and Theory\else ทฤษฎีที่เกี่ยวข้อง\fi}

การทำโครงงาน เริ่มต้นด้วยการศึกษาค้นคว้า ทฤษฎีที่เกี่ยวข้อง หรือ งานวิจัย/โครงงาน ที่เคยมีผู้นำเสนอไว้แล้ว เพื่อให้สามารถนำความรู้ที่ได้จากการศึกษาค้นคว้าเหล่านั้นมาประยุกต์ใช้ในการทำโครงงานได้อย่างมีประสิทธิภาพ และเพื่อให้ผู้อ่าน 
เข้าใจเนื้อหาในบทถัดไปได้ง่ายขึ้น


\section{เซลล์แสงอาทิตย์ (SolarCell) }
เซลล์แสงอาทิตย์ (Solar Cell) หรือ Photovoltaic (PV) Cell 
เป็นอุปกรณ์อิเล็กทรอนิกส์ที่สามารถเปลี่ยนพลังงานแสงอาทิตย์ให้กลายเป็นพลังงานไฟฟ้าผ่านกระบวนการที่ 
เรียกว่า ผลโฟโตโวลเทอิก (Photovoltaic Effect) โดยเซลล์แสงอาทิตย์ประกอบด้วยวัสดุกึ่งตัวนำ 
(Semiconductor) เช่น ซิลิกอน (Silicon) ซึ่งสามารถสร้างกระแสไฟฟ้าได้เมื่อได้รับพลังงานจากโฟตอน 
ของแสงอาทิตย์ โดยระบบพลังงานแสงอาทิตย์นั้นสามารถแบ่งออกได้เป็น 3 ระบบใหญ่ๆ คือ 
  \begin{enumerate}
    \item ระบบออนกริด (On-Grid) : เชื่อมต่อระบบจาก Solarcell เข้ากับ ระบบจากการไฟฟ้า 
  แต่จะไม่สามารถใช้งานได้เมื่อไฟฟ้าดับ และจะต้องมีการขออนุญาตก่อนติดตั้งระบบนี้เสมอ
    \item ระบบออฟกริด (Off-Grid) : ใช้ไฟฟ้าจากระบบ Solarcell อย่างเดียว โดยส่วนมากมักจะมี 
  การเก็บไฟลงในแบตเตอรี่ เพื่อใช้งานในตอนกลางคืน 
    \item ระบบไฮบริด (Hybrid) : เป็นระบบที่ผสมทั้งแบบ ระบบออนกริด (On-Grid),  
  ระบบออฟกริด (Off-Grid) ระบบนี้ สามารถใช้ได้เมื่อไฟดับ และมีระบบแบตเตอรี่ไว้เก็บไฟ แต่มีราคาที่สูง 
  \end{enumerate}
ระบบทั้งสามรูปแบบถูกออกแบบให้มีความเหมาะสมที่แตกต่างกัน ตามลักษณะโครงสร้างพื้นฐานของพื้นที่ติดตั้ง 
และระดับความต้องการในการพึ่งพาพลังงานทดแทนของแต่ละองค์กร




\section{คาร์บอนเครดิต (Carbon Credit)}
คาร์บอนเครดิต (Carbon Credit) เป็นกลไกทางเศรษฐกิจและสิ่งแวดล้อม 
ที่ช่วยลดการปล่อยก๊าซคาร์บอนไดออกไซด์ (CO₂) และก๊าซเรือนกระจกอื่น ๆ ซึ่งเป็นปัจจัยหลักของ 
ภาวะโลกร้อน ระบบนี้กำหนดปริมาณการปล่อยก๊าซของแต่ละองค์กรหรือประเทศตามมาตรฐานที่กำหนด 
หากองค์กรสามารถลดการปล่อยก๊าซได้ต่ำกว่ากำหนด จะได้รับคาร์บอนเครดิตซึ่งสามารถ ขายให้แก่องค์กร 
ที่ปล่อยก๊าซเกินมาตรฐาน การสร้างคาร์บอนเครดิตสามารถดำเนินการผ่านโครงการต่าง ๆ ได้แก่ 
\begin{enumerate}
  \item โครงการพลังงานหมุนเวียน เช่น พลังงานแสงอาทิตย์ พลังงานลม และพลังงานชีวมวล 
ซึ่งช่วยลดการพึ่งพาเชื้อเพลิงฟอสซิล
  \item โครงการด้านป่าไม้และการเกษตร เช่น การปลูกป่า การอนุรักษ์ป่าไม้ 
และเทคนิคทางการเกษตรที่ช่วยกักเก็บคาร์บอนในดิน
  \item อุตสาหกรรมที่ลดการปล่อยมลพิษ 
โดยการพัฒนาเทคโนโลยีและกระบวนการผลิตที่ลดการปล่อยก๊าซเรือนกระจก เช่น การใช้พลังงานสะอาด
  \item โครงการกักเก็บคาร์บอน (Carbon Capture and Storage: CCS) 
ซึ่งดักจับและกักเก็บคาร์บอนจากกระบวนการผลิตหรือลดปริมาณคาร์บอนในบรรยากาศ 
  \item โครงการขององค์กรไม่แสวงหาผลกำไรและชุมชน 
ที่มุ่งเน้นการใช้พลังงานอย่างยั่งยืนและลดผลกระทบต่อสิ่งแวดล้อม
\end{enumerate}
คาร์บอนเครดิตจึงถือเป็นกลไกสำคัญในการขับเคลื่อนนโยบายการลดการปล่อยก๊าซเรือนกระจกสู่ภาคปฏิบัติ 
ตลอดจนเป็นเครื่องมือทางการเงินที่สนับสนุนการพัฒนาอย่างยั่งยืนในระดับสากล




\section{ ความสัมพันธ์ระหว่างสภาพอากาศกับปริมาณการผลิตไฟฟ้าของแผงโซลาร์เซลล์ }
การผลิตไฟฟ้าจากแผงโซลาร์เซลล์ขึ้นอยู่กับปริมาณแสงอาทิตย์ที่สามารถแปลงเป็นพลังงานไฟฟ้าได้ 
อย่างไรก็ตาม ปริมาณพลังงานที่ผลิตได้จริงนั้นได้รับอิทธิพลจากหลายปัจจัยหลัก ได้แก่ 
\begin{enumerate}
  \item ความเข้มของแสง (Light Intensity) แสงแดดที่มีความเข้มสูงช่วยให้แผงโซลาร์เซลล์ผลิตไฟฟ้าได้มากขึ้น 
อย่างไรก็ตาม หากสภาพอากาศมี เมฆปกคลุม หมอก หรือควัน จะลดปริมาณแสงที่เข้าสู่แผง 
ส่งผลให้ประสิทธิภาพการผลิตไฟฟ้าลดลง
  \item อุณหภูมิ (Temperature) อุณหภูมิที่สูงเกินไปอาจส่งผลเสียต่อประสิทธิภาพของแผงโซลาร์เซลล์ โดยทั่วไป 
หากอุณหภูมิของแผงสูงเกิน 25$^\circ$C ประสิทธิภาพการผลิตไฟฟ้าจะลดลง แม้ว่าจะมีแดดจ้า 
แต่หากอุณหภูมิสูงเกินไป ก็อาจส่งผลให้การผลิตพลังงานลดลงมากกว่าที่คาดหวัง
  \item ความชื้นสัมพัทธ์ (Relative Humidity): ความชื้นที่สูงเกินไปการลดทอนความเข้มรังสีดวงอาทิตย์จากการกระเจิงของแสงในชั้นบรรยากาศ และทำให้การระบายความร้อนของแผงโซลาร์เซลล์ลดลง
  ซึ่งนำไปสู่การเพิ่มขึ้นของอุณหภูมิสะสมในเซลล์แสงอาทิตย์
  \item มุมของแสงอาทิตย์ (Tilt and Azimuth Angle) ประสิทธิภาพการผลิตไฟฟ้าของแผงโซลาร์เซลล์จะสูงสุดเมื่อแสงอาทิตย์ตกกระทบ ตั้งฉาก (90 องศา) กับหน้าแผง เนื่องจากความเข้มของพลังงานต่อหน่วยพื้นที่ (Flux Density) จะมีความหนาแน่นสูงสุด ตามหลักการของ 
  กฎกำลังสองผกผันและกฎของคอไซน์ (Lambert's Cosine Law)
  \item ความเร็วลม (Wind Speed): ช่วยระบายความร้อนให้แผงโซลาร์เซลล์ (Cooling Effect) ทำให้แผงทำงานได้ในอุณหภูมิที่ต่ำลงและรักษาประสิทธิภาพไว้ได้ดีขึ้น
\end{enumerate}



\section{เซนเซอร์ (Sensor) และ ไมโครคอนโทรลเลอร์ (Microcontroller)}
เซนเซอร์ (Sensor) คืออุปกรณ์ซึ่งทำหน้าที่สำหรับตรวจจับปริมาณ และสิ่งต่าง ๆ ที่มีการเปลี่ยนแปลงในสภาพแวดล้อมแล้วแปลงเป็นสัญญาณไฟฟ้า องค์ประกอบหลักของเซนเซอร์คือการตรวจจับและวัดสิ่งต่าง ๆ เช่น อุณหภูมิ ความร้อน แสง สี แรงดัน การเคลื่อนที่ ความกว้าง และอื่น ๆ ภายในสิ่งแวดล้อม แล้วจากนั้นเซนเซอร์ก็ทำหน้าที่แปลงสิ่งที่ตรวจจับได้เป็นสัญญาณไฟฟ้าเพื่อนำไปอ่าน วัด ประมวลผล หรือส่งสัญญาณต่อไปยังระบบหรืออุปกรณ์อื่น ๆ เพื่อแสดงเป็นผลลัพธ์ที่ต้องการ
\begin{enumerate}
\item \textbf{BME280}
เซนเซอร์ BME280 เป็นอุปกรณ์ตรวจวัดสภาพแวดล้อมที่สามารถวัดค่าอุณหภูมิ ความชื้นสัมพัทธ์ และความกดอากาศในตัวเดียวกัน โดยใช้หลักการตรวจวัดทางกายภาพร่วมกับวงจรประมวลผลภายในเพื่อแปลงสัญญาณให้เป็นข้อมูลดิจิทัล
สำหรับการวัดอุณหภูมิ เซนเซอร์ใช้หลักการเปลี่ยนแปลงของค่าความต้านทานหรือแรงดันไฟฟ้าตามอุณหภูมิ ส่วนการวัดความชื้นใช้เซนเซอร์แบบ Capacitive ที่วัดการเปลี่ยนแปลงของค่าความจุไฟฟ้าเมื่อความชื้นในอากาศเปลี่ยนแปลง และการวัดความกดอากาศใช้หลักการเปลี่ยนแปลงของแรงดันที่กระทำต่อเซนเซอร์ (Pressure Sensor)
เซนเซอร์ BME280 สามารถสื่อสารข้อมูลผ่านโปรโตคอล I2C ซึ่งเหมาะสำหรับการใช้งานร่วมกับไมโครคอนโทรลเลอร์ เช่น Arduino, ESP8266, ESP32 และ Raspberry Pi
ในด้านคุณสมบัติ เซนเซอร์สามารถวัดอุณหภูมิในช่วง -40 ถึง 85 องศาเซลเซียส ความกดอากาศในช่วง 300 ถึง 1100 hPa และมีความแม่นยำของอุณหภูมิประมาณ $\pm$0.5 องศาเซลเซียส โดยใช้แรงดันไฟฟ้าในการทำงานที่ 3.3 โวลต์ และมีการใช้พลังงานต่ำในโหมด Sleep

\item \textbf{VEML7700}
เซนเซอร์ VEML7700 เป็นอุปกรณ์สำหรับวัดความเข้มแสงในหน่วยลักซ์ (Lux) โดยใช้เทคโนโลยี Filtron ซึ่งถูกออกแบบมาเพื่อปรับการตอบสนองเชิงสเปกตรัมของเซนเซอร์ให้ใกล้เคียงกับการรับรู้แสงของสายตามนุษย์มากที่สุด
หลักการทำงานของเซนเซอร์เริ่มจากโฟโต้ไดโอด (Photodiode) ที่ทำหน้าที่ตรวจจับพลังงานแสงและแปลงเป็นกระแสไฟฟ้า จากนั้นสัญญาณอนาล็อกจะถูกแปลงเป็นสัญญาณดิจิทัลด้วยตัวแปลงสัญญาณอนาล็อกเป็นดิจิทัล (Analog-to-Digital Converter: ADC) ภายในที่มีความละเอียด 16 บิต ซึ่งช่วยให้สามารถวัดค่าความเข้มแสงได้อย่างละเอียดและแม่นยำ รองรับการปรับค่าความไว (Gain) และเวลารวมสัญญาณ (Integration Time) เพื่อให้เหมาะสมกับสภาพแสงที่แตกต่างกัน ตั้งแต่สภาพแสงน้อยมากจนถึงแสงแดดที่มีความเข้มสูง โดยสามารถวัดค่าความเข้มแสงได้ในช่วงประมาณ 0.0036 ถึง 120,000 ลักซ์
การสื่อสารข้อมูลของเซนเซอร์ใช้โปรโตคอล I2C ซึ่งสามารถเชื่อมต่อกับไมโครคอนโทรลเลอร์ เช่น Arduino, ESP8266 หรือ ESP32 และมีฟีเจอร์คำนวณค่าแสงโดยตรง (Lux)

\item \textbf{ความเร็วลม (Wind Speed)}
ความเร็วลม (Wind Speed) คืออัตราเร็วของการเคลื่อนที่ของอากาศจากบริเวณที่มีความกดอากาศสูงไปยังบริเวณที่มีความกดอากาศต่ำ ซึ่งเป็นปัจจัยสำคัญที่มีผลต่อสภาพอากาศและประสิทธิภาพของระบบพลังงานแสงอาทิตย์
การวัดความเร็วลมสามารถทำได้โดยใช้อุปกรณ์ที่เรียกว่า แอนนิโมมิเตอร์ (Anemometer) โดยทั่วไปนิยมใช้แบบถ้วยหมุน (Cup Anemometer) ซึ่งประกอบด้วยถ้วยรับลมติดตั้งอยู่กับแกนหมุน เมื่อมีลมพัดมากระทบ ถ้วยจะหมุนตามแรงลม โดยความเร็วในการหมุนจะเป็นสัดส่วนโดยตรงกับความเร็วลม สัญญาณการหมุนของแกนจะถูกแปลงเป็นสัญญาณไฟฟ้า เช่น สัญญาณพัลส์ (Pulse Signal) หรือแรงดันไฟฟ้า ซึ่งสามารถนำไปคำนวณเป็นค่าความเร็วลมในหน่วยเมตรต่อวินาที (m/s) ได้

\item \textbf{ทิศทางลม (Wind Direction)}
ทิศทางลม (Wind Direction) คือทิศทางที่ลมพัดมา เช่น ลมเหนือพัดมาจากทิศเหนือ การวัดทิศทางลมสามารถทำได้โดยใช้อุปกรณ์ที่เรียกว่า ศรลม (Wind Vane) ซึ่งใบพัดจะหันไปตามแรงลม โดยใช้หลักการสมดุลของแรงปะทะที่หางใบพัดกับหัวใบพัด หางใบพัดซึ่งมีพื้นที่มากกว่าจะรับแรงลมมากกว่า ทำให้เกิดสมดุลของแรงและส่งผลให้หัวศรหันไปในทิศทางที่ลมพัดมา ตำแหน่งเชิงมุมของศรลมจะถูกแปลงเป็นสัญญาณไฟฟ้า เช่น แรงดันไฟฟ้าแบบแอนะล็อก หรือสัญญาณดิจิทัล จากนั้นสามารถนำไปคำนวณเป็นค่าทิศทางลมในรูปแบบองศา (0–360 องศา) หรือทิศหลัก เช่น เหนือ ใต้ ตะวันออก และตะวันตก

\item \textbf{PT100}
PT100 เป็นเซนเซอร์วัดอุณหภูมิประเภท RTD (Resistance Temperature Detector) ที่ทำจากแพลตตินัม โดยอาศัยหลักการที่ว่าค่าความต้านทานไฟฟ้าของวัสดุจะเปลี่ยนแปลงตามอุณหภูมิในลักษณะเชิงเส้น กล่าวคือ เมื่ออุณหภูมิสูงขึ้น ค่าความต้านทานจะเพิ่มขึ้น
เซนเซอร์ PT100 มีค่าความต้านทานมาตรฐานเท่ากับ 100 โอห์ม ที่อุณหภูมิ 0 องศาเซลเซียส โครงสร้างของเซนเซอร์ประกอบด้วยเส้นลวดแพลตตินัมพันรอบแกนเซรามิกหรือแก้ว ซึ่งแพลตตินัมมีคุณสมบัติเด่นในด้านความเสถียรทางเคมีและสามารถทนต่ออุณหภูมิสูงได้ดี
ในการวัดอุณหภูมิ ระบบจะจ่ายกระแสไฟฟ้าปริมาณเล็กน้อยผ่านเซนเซอร์ และวัดแรงดันไฟฟ้าที่เกิดขึ้นเพื่อนำไปคำนวณหาค่าความต้านทานตามกฎของโอห์ม ($R = V/I$) จากนั้นค่าความต้านทานที่ได้จะถูกนำไปเปรียบเทียบกับตารางมาตรฐาน เช่น DIN/IEC 60751 เพื่อแปลงเป็นค่าอุณหภูมิ เซนเซอร์ PT100 สามารถวัดอุณหภูมิได้ในช่วงประมาณ -200 ถึง 850 องศาเซลเซียส และมีความแม่นยำสูง เหมาะสำหรับงานที่ต้องการความเที่ยงตรงในการวัดอุณหภูมิ
\end{enumerate}

ไมโครคอนโทรลเลอร์ (Microcontroller: MCU) คือคอมพิวเตอร์ขนาดเล็กบนชิปเดียว ออกแบบมาเพื่อจัดการงานเฉพาะภายในระบบฝังตัวโดยไม่ต้องใช้ระบบปฏิบัติการที่ซับซ้อน
โดยได้รวม CPU, หน่วยความจำ (RAM/ROM) และอุปกรณ์อินพุต/เอาต์พุต (Input/Output: I/O) ทำหน้าที่หลักในการควบคุมอุปกรณ์อิเล็กทรอนิกส์, รับข้อมูลผ่านเซนเซอร์, ประมวลผล และส่งคำสั่งไปสั่งงานอุปกรณ์ภายนอกตามโปรแกรมที่เขียนไว้ นิยมใช้ในระบบฝังตัว (Embedded System) 
\begin{enumerate}
\item \textbf{ESP32}
ESP32 เป็นไมโครคอนโทรลเลอร์ประสิทธิภาพสูงที่รวมความสามารถในการสื่อสารแบบไร้สาย ได้แก่ Wi-Fi และ Bluetooth (Classic และ BLE) ไว้ภายในชิปเดียว โดยมีสถาปัตยกรรมแบบ 32 บิต ชนิด Dual-Core ซึ่งสามารถทำงานที่ความเร็วสูงสุด 240 MHz และมีหน่วยความจำภายใน (RAM) ขนาด 512 KB ทำให้เหมาะสำหรับงานที่ต้องการการประมวลผลข้อมูลอย่างรวดเร็วและมีประสิทธิภาพสูง โดยเฉพาะในระบบ Internet of Things (IoT)
ESP32 ใช้หน่วยประมวลผล Tensilica Xtensa LX6 ซึ่งประกอบด้วยหน่วยประมวลผล 2 แกน (Dual-Core) ทำให้สามารถแบ่งการทำงานออกเป็นหลายส่วนได้อย่างมีประสิทธิภาพ เช่น การประมวลผลข้อมูลจากเซ็นเซอร์ และการสื่อสารข้อมูลผ่านเครือข่ายไร้สายพร้อมกัน ช่วยลดความล่าช้าและเพิ่มความเสถียรของระบบโดยรวม
ในด้านการสื่อสาร ESP32 สามารถทำงานได้ทั้งในโหมด Access Point (AP) สำหรับสร้างเครือข่าย Wi-Fi และโหมด Station (STA) สำหรับเชื่อมต่อกับเครือข่าย Wi-Fi ภายนอก อีกทั้งยังรองรับโปรโตคอลที่นิยมใช้ในระบบ IoT เช่น MQTT และ HTTP สำหรับการส่งข้อมูลไปยังเซิร์ฟเวอร์ นอกจากนี้ ESP32 ยังรองรับระบบปฏิบัติการแบบเรียลไทม์ (FreeRTOS) ซึ่งช่วยในการจัดการลำดับความสำคัญของงาน (Task Scheduling) ทำให้สามารถทำงานหลายกระบวนการพร้อมกันได้อย่างมีประสิทธิภาพ
ในส่วนของฮาร์ดแวร์ ESP32 มีขาเชื่อมต่อ (GPIO) มากกว่า 30 ขา และรองรับการทำงานที่หลากหลาย เช่น ADC สำหรับอ่านค่าสัญญาณอนาล็อก, DAC สำหรับสร้างสัญญาณอนาล็อก, PWM สำหรับควบคุมอุปกรณ์ รวมถึงรองรับโปรโตคอลการสื่อสารต่าง ๆ เช่น UART, I2C และ SPI นอกจากนี้ ESP32 ยังมีโหมดประหยัดพลังงาน เช่น Deep Sleep ซึ่งใช้กระแสไฟฟ้าต่ำมาก (ประมาณ 2.5 ไมโครแอมแปร์) ทำให้เหมาะสำหรับระบบที่ใช้พลังงานจากแบตเตอรี่ โดยตัวอุปกรณ์ทำงานที่แรงดันไฟฟ้า 3.3 โวลต์ ซึ่งต่ำกว่าไมโครคอนโทรลเลอร์ทั่วไป ส่งผลให้มีประสิทธิภาพด้านการใช้พลังงานที่ดี
\end{enumerate}


\section{ระบบอินเทอร์เน็ตของสรรพสิ่ง (Internet of Things: IoT)}
IoT หรืออินเทอร์เน็ตในทุกสิ่ง (Internet of Things) หมายถึงเครือข่ายรวมของอุปกรณ์ที่เชื่อมต่อถึงกันและเทคโนโลยีที่อำนวยความสะดวกในการสื่อสารระหว่างอุปกรณ์กับระบบคลาวด์ ตลอดจนระหว่างอุปกรณ์ด้วยกันเอง
โครงสร้างของระบบ IoT ในงานวิจัยนี้ประกอบด้วย 3 ส่วนหลัก ได้แก่ ส่วนอุปกรณ์ตรวจวัด (Sensors) ส่วนประมวลผลเบื้องต้น (Microcontroller) และส่วนเซิร์ฟเวอร์ (Server) โดยเซนเซอร์จะทำหน้าที่ตรวจวัดข้อมูลสภาพอากาศ เช่น อุณหภูมิ ความชื้น ความกดอากาศ ความเข้มแสง ความเร็วลม ทิศทางลม และอุณหภูมิของแผงโซลาร์เซลล์ จากนั้นข้อมูลจะถูกส่งไปยังไมโครคอนโทรลเลอร์เพื่อทำการรวบรวมและจัดรูปแบบข้อมูลก่อนส่งผ่านเครือข่ายอินเทอร์เน็ต
การส่งข้อมูลในระบบสามารถดำเนินการผ่านเครือข่ายไร้สาย (Cellular) โดยใช้โปรโตคอล เช่น HTTPS เพื่อส่งข้อมูลไปยังเซิร์ฟเวอร์กลาง ซึ่งทำหน้าที่จัดเก็บข้อมูลในฐานข้อมูลและนำข้อมูลไปใช้ในการพยากรณ์ปริมาณการผลิตไฟฟ้าจากแผงโซลาร์เซลล์ด้วยโมเดล Machine Learning
กระบวนการไหลของข้อมูล (Data Flow) ในระบบเริ่มต้นจากการตรวจวัดข้อมูลโดยเซนเซอร์ จากนั้นข้อมูลจะถูกส่งไปยังไมโครคอนโทรลเลอร์เพื่อประมวลผลเบื้องต้น และส่งต่อไปยังเซิร์ฟเวอร์เพื่อจัดเก็บและวิเคราะห์ ผลลัพธ์ที่ได้จากการพยากรณ์สามารถนำไปแสดงผลผ่านระบบแสดงผล (Dashboard) เพื่อช่วยให้ผู้ใช้งานสามารถติดตามและวางแผนการใช้พลังงานได้อย่างมีประสิทธิภาพ
ดังนั้น ระบบ IoT จึงเป็นส่วนสำคัญที่ทำให้สามารถเชื่อมโยงข้อมูลจากโลกจริงเข้าสู่ระบบดิจิทัล 
และสนับสนุนการทำงานของระบบพยากรณ์และการบริหารจัดการพลังงานได้อย่างมีประสิทธิภาพ




\section{ประเภทของข้อมูลเวลา}
การแบ่งประเภทของข้อมูลโดยอาศัยระยะเวลาที่จัดเก็บข้อมูลจะสามารถแบ่งออกได้เป็น 2 ประเภท ได้แก่ ข้อมูลอนุกรมเวลา และข้อมูลตัดขวาง
\begin{enumerate}
  \item ข้อมูลอนุกรมเวลา (time series data)
คือ ข้อมูลที่ได้จากการจัดเก็บตามลำดับเวลาต่อเนื่องกันเป็นช่วงๆ เช่น ทุก 7 วัน ทุก 1 เดือน ทุก 1 ปี ฯลฯ ข้อมูลประเภทนี้ช่วยแสดงให้เห็นการเปลี่ยนแปลงของข้อมูล
เช่น การเพิ่มขึ้นของข้อมูล การลดลงของข้อมูล ค่าสูงสุด ค่าต่ำสุด และแนวโน้มของข้อมูลในช่วงเวลาที่กำหนด
  \item ข้อมูลตัดขวาง (cross-sectional data) ข้อมูลที่บอกลักษณะ สถานะ หรือ สภาพของสิ่งที่สนใจ ณ เวลาใดเวลาหนึ่ง เช่น ข้อมูลอุณหภูมิ ของประเทศไทย ณ วันที่ 1 มกราคม 2025 ข้อมูลประเภทนี้ช่วยให้เห็นภาพรวมของข้อมูลในช่วงเวลาที่กำหนด และสามารถนำไปเปรียบเทียบกับข้อมูลในช่วงเวลาที่ต่างกันได้
\end{enumerate}
ข้อมูลอนุกรมเวลามักมีลักษณะสำคัญ เช่น แนวโน้ม (Trend) และฤดูกาล (Seasonality) 
ซึ่งเป็นปัจจัยสำคัญในการสร้างแบบจำลองพยากรณ์

\section{\ifenglish Recurrent Neural Network (RNN)\else โครงข่ายประสาทเทียมแบบย้อนกลับ (Recurrent Neural Network)\fi}
โครงข่ายประสาทเทียมแบบย้อนกลับ (Recurrent Neural Network: RNN) เป็นประเภทหนึ่งของโครงข่ายประสาทเทียมที่ถูกออกแบบมาเพื่อประมวลผลข้อมูลที่มีลักษณะเป็นลำดับ (Sequential Data) หรือข้อมูลอนุกรมเวลา (Time-series Data) โดยจุดเด่นที่แตกต่างจากโครงข่ายประสาทเทียมแบบส่งต่อ (Feedforward Neural Network) ทั่วไปคือ RNN มีส่วนของหน่วยความจำ (Memory) ที่ทำหน้าที่เก็บข้อมูลจากลำดับก่อนหน้าเพื่อนำมาใช้ร่วมกับการประมวลผลในลำดับปัจจุบัน
ในการคำนวณแต่ละขั้นตอนของเวลา (Time Step: $t$) โมเดลจะรับข้อมูลนำเข้า ($x_t$) ร่วมกับสถานะที่ถูกซ่อนอยู่ (Hidden State) จากขั้นตอนก่อนหน้า ($h_{t-1}$) เพื่อสร้างสถานะที่ซ่อน (Hidden State) ใหม่ ($h_t$) สำหรับส่งต่อไปยังลำดับถัดไป โดยสามารถเขียนแทนด้วยสมการทางคณิตศาสตร์ได้ดังนี้:
\begin{equation}
    h_t = f(W_h h_{t-1} + W_x x_t + b)
\end{equation}
\noindent โดยที่รายละเอียดของตัวแปรมีดังนี้:
\begin{itemize}
    \item $h_t$ คือ สถานะซ่อนเร้น (Hidden State) ณ เวลา $t$
    \item $x_t$ คือ ข้อมูลนำเข้า (Input Vector) ณ เวลา $t$
    \item $h_{t-1}$ คือ สถานะซ่อนเร้นจากลำดับก่อนหน้า (Previous Hidden State)
    \item $W_h, W_x$ คือ เมทริกซ์น้ำหนัก (Weight Matrices) ของโมเดล
    \item $b$ คือ ค่าอคติ (Bias Vector)
    \item $f$ คือ ฟังก์ชันกระตุ้น (Activation Function) เช่น $\tanh$ หรือ $\text{ReLU}$
\end{itemize}
ด้วยสถาปัตยกรรมดังกล่าว ทำให้ RNN มีความเหมาะสมอย่างยิ่งสำหรับงานพยากรณ์ปริมาณการผลิตไฟฟ้าจากเซลล์แสงอาทิตย์ เนื่องจากสามารถเรียนรู้ความสัมพันธ์ของปัจจัยสภาพอากาศที่มีความต่อเนื่องทางเวลาได้




\section{ตัวชี้วัดประสิทธิภาพของแบบจำลอง (Model Evaluation Metrics) }
ในการประเมินประสิทธิภาพของแบบจำลองการพยากรณ์ปริมาณการผลิตไฟฟ้าจากเซลล์แสงอาทิตย์ โครงงานนี้เลือกใช้ตัวชี้วัดมาตรฐานทางสถิติเพื่อวัดความคลาดเคลื่อนและความแม่นยำของผลการทำนาย ดังนี้
\begin{enumerate}
  \item ค่าความคลาดเคลื่อนสมบูรณ์เฉลี่ย (Mean Absolute Error: MAE) 
  คือ ตัวชี้วัดที่ใช้สำหรับวัดค่าเฉลี่ยของขนาดความคลาดเคลื่อนระหว่างค่าที่พยากรณ์ได้กับค่าจริงที่เกิดขึ้น โดยการคำนวณจะนำค่าสัมบูรณ์ของผลต่างในแต่ละข้อมูลมาหาค่าเฉลี่ย ซึ่งค่า MAE ที่ต่ำหมายถึงแบบจำลองมีความแม่นยำสูง โดยมีสมการคำนวณดังนี้
  \begin{equation}
  \text{MAE} = \frac{1}{n} \sum_{i=1}^{n} |y_i - \hat{y}_i|
  \end{equation}
  \noindent โดยที่:
  \begin{itemize}
  \item $n$ คือ จำนวนข้อมูลทั้งหมดที่ใช้ในการทดสอบ
  \item $y_i$ คือ ค่าจริง (Actual Value) ของลำดับที่ $i$
  \item $\hat{y}_i$ คือ ค่าที่ได้จากการพยากรณ์ (Predicted Value) ของลำดับที่ $i$
  \end{itemize}

  \item สัมประสิทธิ์การกำหนด (coefficient of determination) หรือ R-squared เป็นมาตรวัดทางสถิติที่ใช้ในแมชชีนเลิร์นนิง (Machine Learning) สำหรับประเมินคุณภาพของแบบจำลองการถดถอย ใช้ในการวัดว่าแบบจำลองนั้นเหมาะสมกับข้อมูลได้ดีเพียงใด และสามารถทำนายผลลัพธ์ในอนาคตได้ดีเพียงใด
  ค่า R-squared อยู่ในช่วง 0 ถึง 1 โดยยิ่งค่า R-squared ใกล้เคียงกับหนึ่งมากเท่าไหร่ แบบจำลองก็จะยิ่งเหมาะสมกับข้อมูลได้ดี
  ค่า R-squared คำนวณทางคณิตศาสตร์โดยการเปรียบเทียบผลรวมกำลังสองของความคลาดเคลื่อน (Sum of Squares Error: SSE) กับผลรวมกำลังสองทั้งหมด (Total Sum of Squares: SST)
  \begin{equation}R^2 = 1 - \frac{\sum (y_i - \hat{y}_i)^2}{\sum (y_i - \bar{y})^2}
  \end{equation}
  \noindent โดยที่:
  \begin{itemize}\item $\bar{y}$ คือ ค่าเฉลี่ยของข้อมูลจริงทั้งหมด (Mean of Actual Values)
    \item $\sum (y_i - \hat{y}_i)^2$ คือ ส่วนเบี่ยงเบนกำลังสองของความคลาดเคลื่อน (SSE)
    \item $\sum (y_i - \bar{y})^2$ คือ ส่วนเบี่ยงเบนกำลังสองทั้งหมด (SST)
  \end{itemize}

  \end{enumerate}
